\section{An order trial with prescribed success probability}
\label{sec:order}

For a unit $a$ modulo $m$, write $\ord_m(a)$ for its multiplicative order.
We construct a trial whose probability of returning the true order is an
explicit dyadic function of that order.  This additional control of the
probability will allow us to make the final factoring algorithm exact.
As usual, let
\[
 \varphi(d)=\bigl|\{0\le j<d:\gcd(j,d)=1\}\bigr|,
 \qquad \varphi(1)=1.
\]

\begin{proposition}\label{prop:order-trial}
Let $2\le m<B$ and let $a$ be a unit modulo $m$, of order $r$.
There is a uniform quantum trial of worst-case cost polynomial in $n$
that returns either no candidate or an integer $d$ satisfying
$1\le d<B$ and $r\mid d$.  Moreover,
\[
 \Pr(d=r)=\varphi(r)\lambda(r),
 \qquad
 \lambda(d)=2^{-\lceil\log_2 d\rceil-10}.
\]
The trial uses only the fixed finite gate set specified above and does not
require $r$ or any factorization as input.
\end{proposition}

The construction has two parts.  First we replace the Fourier phases in
order sampling by a piecewise linear phase function.  The resulting
probabilities admit exact dyadic arithmetic.  We then add two small
classical sampling branches to make the mass at each reduced fraction of
denominator $r$ exactly $\lambda(r)$.  The proof of
Proposition~\ref{prop:order-trial} is completed at the end of this section.

\subsection{A phase function implemented by a circuit}

Set
\[
 Q=B^{16}=2^b,\qquad b=16n.
\]
Let $g:\R\to\C$ be the continuous, $1$-periodic function obtained by
linear interpolation of $\exp(-2\pi\mathrm i z)$ at the quarter-integers.
Explicitly, for $q\in\Z$ and $0\le t\le1$,
\begin{equation}\label{eq:linear-phase}
 g\bigl((q+t)/4\bigr)
   =(1-t)(-\mathrm i)^q+t(-\mathrm i)^{q+1}.
\end{equation}
In particular, $g(z+1/2)=-g(z)$.

The controlled-shift mechanism is a form of phase kickback~\cite[Section~7]{CEMM}.
Here the auxiliary state is not a shift eigenstate: its overlap with a
shifted copy will give $g$. The success test in the next lemma is therefore
part of the matrix's probability, rather than an instruction to apply an
ideal Fourier transform.

\begin{lemma}\label{lem:transform}
There is a circuit of size polynomial in $b$ on a $b$-qubit data register
and ancillary registers initially in the all-zero state, with a designated
ancillary success outcome, such
that its unnormalized matrix on that outcome is
\[
 A_{k,x}=Q^{-1/2}g(kx/Q),\qquad 0\le x,k<Q.
\]
The circuit uses only the fixed finite gate set.
\end{lemma}

\begin{proof}
Prepare a $b$-qubit phase register in the state
\[
 \ket{\chi}=Q^{-1/2}\sum_{y=0}^{Q-1}
                  \mathrm i^{\lfloor4y/Q\rfloor}\ket{y}.
\]
This requires Hadamards on its bits, followed by a $Z=S^2$ phase on the
most significant bit and an $S$ phase on the next bit.
For $D\in\Z$, let $T_D$ denote the forward shift
$\ket y\mapsto\ket{y+D\bmod Q}$ on this register.  We first verify that
\begin{equation}\label{eq:phase-overlap}
 \bra{\chi}T_D\ket{\chi}=g(D/Q).
\end{equation}
Write $D=qQ/4+s$ modulo $Q$, where $0\le q<4$ and $0\le s<Q/4$.
Among the $Q$ positions, a fraction $1-4s/Q$ crosses exactly $q$
quarter-bins, and the remaining fraction $4s/Q$ crosses $q+1$.
The corresponding products of the initial phase with the conjugate final
phase are $\mathrm i^{-q}$ and $\mathrm i^{-q-1}$, respectively.
Their average is precisely~\eqref{eq:linear-phase} at $D/Q$.
The sign is fixed by the use of a forward shift: a shift by $Q/4$ has
overlap $-\mathrm i$.

Number input and output bits from the units bit:
$x=\sum_{i=0}^{b-1}2^ix_i$ and $k=\sum_{j=0}^{b-1}2^jk_j$.
For $j=0,\ldots,b-1$, perform the following operations in order:
apply a Hadamard to the physical bit that initially carried $x_{b-1-j}$,
regard its new value as $k_j$, and, for every $i<b-1-j$, apply the shift
\[
 T_{2^{i+j}x_i k_j}
\]
to the phase register.  Both controls are available: $k_j$ is the bit
just produced, whereas $x_i$ has not yet been overwritten.  The logical
output bits are thus in the reverse physical order from the input bits;
a wire reversal, or ordinary SWAP gates, puts them in the conventional
order.

For fixed input $x$ and output $k$, each data bit has exactly its specified
value before and after its single Hadamard.  The product of the Hadamard
matrix elements is $Q^{-1/2}(-1)^h$, and the total phase-register shift is
$D$, where
\[
 h=\sum_{i+j=b-1}x_i k_j,
 \qquad
 D=\sum_{i+j<b-1}2^{i+j}x_i k_j.
\]
All omitted terms in the product $xk$ are accounted for by
\[
 xk\equiv D+(Q/2)h\pmod Q.
\]
Using~\eqref{eq:phase-overlap} and $g(z+1/2)=-g(z)$ gives the claimed
matrix entry after projecting the phase register back onto $\ket\chi$:
\[
 Q^{-1/2}(-1)^h\bra\chi T_D\ket\chi
       =Q^{-1/2}g(xk/Q).
\]
To perform that test, undo the preparation of $\ket\chi$ and designate
the all-zero phase-register outcome as success.  Its actual probability
is retained; there is no instruction to enforce that outcome.

There are $O(b^2)$ controlled shifts.  Each is a reversible modular
binary addition on $b$ bits, with its scratch bits uncomputed afterwards,
and therefore has polynomial cost in the stated finite gate set.  All
other operations also have polynomial cost.
\end{proof}

\subsection{Recovering fraction labels and their probabilities}

For every reduced fraction $j/d$ with $0\le j<d<B$, define
\begin{equation}\label{eq:order-bins}
 k(d,j)=\left\lfloor\frac{jQ}{d}+\frac12\right\rfloor,
 \qquad
 \eta=k(d,j)-\frac{jQ}{d}.
\end{equation}
Thus $|\eta|\le1/2$, and $0\le k(d,j)<Q$.
We call $(d,j)$ the label of this bin.

\begin{lemma}\label{lem:bins}
Given $0\le k<Q$, one can determine in polynomial bit time whether
$k=k(d,j)$ for a reduced fraction $0\le j<d<B$, and recover $(d,j)$
when it exists.  Such a label is unique.
\end{lemma}

\begin{proof}
Compute the finite simple continued fraction of $k/Q$, and let $p/v$ be
its last convergent with denominator less than $B$.  If all convergents
have denominator less than $B$, use the terminal one.  We have
\begin{equation}\label{eq:convergent-bound}
 \left|\frac{k}{Q}-\frac pv\right|\le\frac1{vB}.
\end{equation}
Indeed, suppose $p/v$ is not terminal.  Let $(p',v')$ be the preceding
numerator--denominator pair, using the formal pair $(1,0)$ when $p/v$
is the first convergent.  The remaining continued-fraction tail $\theta$
gives
\[
 \frac{k}{Q}=\frac{p\theta+p'}{v\theta+v'},
 \qquad |pv'-p'v|=1.
\]
The next denominator is $v c+v'\ge B$, where $c$ is the integer part of
$\theta$.  Consequently the error is
$1/[v(v\theta+v')]\le1/[v(vc+v')]\le1/(vB)$.
The determinant identity follows inductively from the convergent
recurrence, which replaces each numerator--denominator pair by $c$
times the latest pair plus the preceding pair.  If the convergent is
terminal, its error is zero.

Suppose $k$ has a label $(d,j)$.  If $j/d\ne p/v$, integrality of the
cross-product gives
\[
 \left|\frac jd-\frac pv\right|
 \ge\frac1{dv}\ge\frac1{(B-1)v}
 >\frac1{vB}+\frac1{2Q}.
\]
For the strict inequality, note that
$2Q>vB(B-1)$ because $v<B$ and $Q=B^{16}$.
This contradicts~\eqref{eq:convergent-bound} and
$|k/Q-j/d|\le1/(2Q)$.  Thus $j/d=p/v$.
It remains only to check $0\le p<v<B$ and the exact rule
\eqref{eq:order-bins}; these checks accept precisely the valid labels.
They also prove uniqueness.  When $k=0$, the terminal convergent is
$0/1$, so the case $d=1$ is included.
The Euclidean algorithm and all these checks use integers with $O(n)$
bits and take polynomial bit time.
\end{proof}

We now apply Lemma~\ref{lem:transform} to the state
\[
 Q^{-1/2}\sum_{x=0}^{Q-1}\ket x\ket{a^x\bmod m}.
\]
Reversible modular powering prepares this state in polynomial time;
its temporary work is uncomputed, leaving no additional information
about $x$.  This is the order-sampling setup of the quantum factoring
algorithm~\cite{Shor}, with the transform of Lemma~\ref{lem:transform}
in place of the Fourier transform.

For any reduced label $(d,j)$ and $0\le t<d$, define the arithmetic
quantity
\begin{equation}\label{eq:order-P}
 P(d,j,t)=\left|\frac1Q
       \sum_{\substack{0\le x<Q\\x\equiv t\pmod d}}
            g\bigl(k(d,j)x/Q\bigr)\right|^2.
\end{equation}
When $d=r=\ord_m(a)$, it has the following direct probabilistic meaning:
\begin{equation}\label{eq:residue-probability}
 \Pr\bigl(\text{phase-register success and }k=k(r,j)\bigr)
       =\sum_{t=0}^{r-1}P(r,j,t).
\end{equation}
To see this, the unnormalized second-register vector for a fixed output
$k$ is
\[
 \frac1Q\sum_{x=0}^{Q-1}g(kx/Q)\ket{a^x\bmod m}.
\]
The labels $a^x\bmod m$ coincide exactly on residue classes modulo $r$;
the distinct classes give orthogonal basis states.  Taking its squared
norm proves~\eqref{eq:residue-probability}, including the actual
probability of the phase-register test.

The definition~\eqref{eq:order-P} will also be used for candidate
values $d\ne r$; for those values no physical-probability interpretation
is asserted.  The next step is to compute and correct the masses in
\eqref{eq:residue-probability} using arithmetic on candidate labels,
without knowing the true order.
